\documentclass[12pt,a4paper]{article}

% ============================================
% PACKAGES
% ============================================
\usepackage[utf8]{inputenc}
\usepackage[T1]{fontenc}
\usepackage{upquote}
\usepackage[english]{babel}
\usepackage{geometry}
\usepackage{graphicx}
\usepackage{amsmath}
\usepackage{amssymb}
\usepackage{listings}
\usepackage{xcolor}
\usepackage{hyperref}
\usepackage{tocloft}
\usepackage{fancyhdr}
\usepackage{titlesec}
\usepackage{caption}
\usepackage{subcaption}
\usepackage{booktabs}
\usepackage{array}
\usepackage{multirow}
\usepackage{float}

% ============================================
% PAGE LAYOUT
% ============================================
\geometry{
    a4paper,
    left=30mm, right=15mm,
    top=20mm,  bottom=20mm,
    headheight=15pt
}

% ============================================
% CODE STYLE
% ============================================
\definecolor{codegreen}{rgb}{0,0.6,0}
\definecolor{codegray}{rgb}{0.5,0.5,0.5}
\definecolor{codepurple}{rgb}{0.58,0,0.82}
\definecolor{backcolour}{rgb}{0.05,0.05,0.05}
\definecolor{codeorange}{rgb}{1,0.5,0}
\definecolor{codeyellow}{rgb}{0.9,0.8,0.2}
\definecolor{codeblue}{rgb}{0.4,0.6,1}

\lstdefinelanguage{nasm}{
    morekeywords={section,global,bits,mov,add,sub,imul,idiv,cdq,cmp,jmp,je,jne,
                  jg,jge,jl,jle,jb,ja,jnz,jz,jnc,push,pop,call,ret,int,xor,
                  inc,dec,test,div,daa,db,dw,dd,equ,resb,resw,resd,
                  byte,word,dword,main},
    sensitive=false,
    morecomment=[l]{;},
    morestring=[b]",
    morestring=[b]',
}

\lstdefinestyle{asmstyle}{
    backgroundcolor=\color{backcolour},
    commentstyle=\color{codegreen},
    keywordstyle=\color{codepurple},
    numberstyle=\tiny\color{codegray},
    stringstyle=\color{codeorange},
    basicstyle=\ttfamily\footnotesize\color{white},
    breaklines=true,
    captionpos=b,
    keepspaces=true,
    numbers=left,
    numbersep=5pt,
    showstringspaces=false,
    tabsize=4,
    language=nasm,
}
\lstset{style=asmstyle}

% ============================================
% SECTION FORMATTING
% ============================================
\titleformat{\section}{\normalfont\Large\bfseries\uppercase}{\thesection}{1em}{}
\titleformat{\subsection}{\normalfont\large\bfseries}{\thesubsection}{1em}{}
\titleformat{\subsubsection}{\normalfont\normalsize\bfseries}{\thesubsubsection}{1em}{}

% ============================================
% HYPERLINKS
% ============================================
\hypersetup{
    colorlinks=true, linkcolor=black,
    urlcolor=cyan,
    pdftitle={Laboratory Work 5 - Assembly and BCD Numbers},
}

% ============================================
% HEADER/FOOTER
% ============================================
\pagestyle{fancy}
\fancyhf{}
\fancyhead[R]{\thepage}
\renewcommand{\headrulewidth}{0pt}

% ============================================
% CAPTION SETUP
% ============================================
\captionsetup[figure]{font=small,labelfont=bf,skip=3pt}

\begin{document}

% ============================================
% TITLE PAGE
% ============================================
\begin{titlepage}
    \centering
    \vspace*{1cm}
    {\Large Ministerul Educa\c{t}iei \c{s}i Cercet\u{a}rii al Republicii Moldova\\}
    \vspace{0.5cm}
    {\Large Universitatea Tehnic\u{a} a Moldovei\\}
    \vspace{0.5cm}
    {\Large Facultatea Calculatoare, Informatic\u{a} \c{s}i Microelectronic\u{a}\\}
    \vspace{3cm}
    {\LARGE \textbf{Laboratory work No. 5:}\\}
    \vspace{0.5cm}
    {\LARGE \textbf{Assembly Arithmetic \&}\\}
    \vspace{0.3cm}
    {\LARGE \textbf{BCD Number Representation}\\}
    \vspace{0.3cm}
    {\LARGE \textbf{in NASM}\\}
    \vfill
    \begin{flushleft}
    \hspace{7cm}\textbf{Elaborated:}\\
    \hspace{7cm}st. gr. FAF-243 Kushnirenko Ecaterina\\
    \vspace{1cm}
    \hspace{7cm}\textbf{Verified:}\\
    \hspace{7cm}asist. univ. Fi\c{s}tic Cristofor\\
    \end{flushleft}
    \vspace{1cm}
    {\large Chi\c{s}in\u{a}u - 2026}
\end{titlepage}

% ============================================
% TABLE OF CONTENTS
% ============================================
\newpage
\tableofcontents
\thispagestyle{empty}

\newpage
\setcounter{page}{3}

% ============================================
\section{INTRODUCTION}
% ============================================

\subsection{Registers and Arithmetic in x86 Assembly}

In x86 Assembly Language, all computation takes place directly in CPU registers --- small,
ultra-fast storage locations built into the processor. Unlike high-level languages that
hide arithmetic behind operators, Assembly requires the programmer to explicitly load
operands into registers, apply an instruction, and store the result. This gives complete
visibility into every step of a computation.

The general-purpose registers used for arithmetic in 32-bit (IA-32) mode are:

\begin{table}[H]
\centering
\caption{General-purpose registers used in this laboratory}
\small
\begin{tabular}{|l|p{10cm}|}
\hline
\textbf{Register} & \textbf{Primary Role} \\
\hline
EAX & Accumulator --- holds operands and results; \texttt{idiv} places quotient here \\
\hline
EBX & General operand storage \\
\hline
ECX & Counter and temporary operand \\
\hline
EDX & Extended accumulator --- \texttt{cdq} sign-extends EAX into EDX:EAX before division;
\texttt{idiv} places remainder here \\
\hline
ESI & Save register --- used to preserve intermediate results across operations \\
\hline
\end{tabular}
\end{table}

\subsection{BCD (Binary Coded Decimal) Representation}

Binary Coded Decimal (BCD) is a class of encodings in which each decimal digit (0--9) is
represented individually by its binary equivalent. The most common variant is the
\textbf{8421 BCD} code, where each digit occupies exactly four bits (a nibble), weighted
$2^3, 2^2, 2^1, 2^0$.

\begin{table}[H]
\centering
\caption{8421 BCD encoding table}
\small
\begin{tabular}{|c|c||c|c|}
\hline
\textbf{Decimal} & \textbf{BCD} & \textbf{Decimal} & \textbf{BCD} \\
\hline
0 & 0000 & 5 & 0101 \\
\hline
1 & 0001 & 6 & 0110 \\
\hline
2 & 0010 & 7 & 0111 \\
\hline
3 & 0011 & 8 & 1000 \\
\hline
4 & 0100 & 9 & 1001 \\
\hline
\end{tabular}
\end{table}

The six bit patterns 1010--1111 are invalid in BCD. When arithmetic produces these
patterns, a correction of \texttt{+6} (0110) is applied to skip over the invalid states
and return to a valid BCD digit.

% ============================================
\section{PURPOSE OF THE REPORT}
% ============================================

This report documents Laboratory Work No.\ 5, Variant 14.  The objectives are:

\begin{itemize}
    \renewcommand{\labelitemi}{--}
    \item Write a NASM program that computes the register expression $A = (A+B)-(C+D)$
    \item Implement three arithmetic expressions selected according to the variant rules
    \item Perform manual BCD conversions, additions, and subtractions using 9's complement
    \item Understand the structure of a BCD adder circuit
    \item Write a NASM program that adds two packed 8-bit BCD numbers using the \texttt{DAA}
    instruction
\end{itemize}

\subsection{Variant Selection}

Variant number: \textbf{14} (even).

\begin{itemize}
    \renewcommand{\labelitemi}{--}
    \item \textbf{Expression 1:} Variant 14 $\Rightarrow$ \textbf{Expression 14}
    \item \textbf{Expression 2:} 14 is even $\Rightarrow$ add 2: $14+2 = 16 > 15$, so
    divide by 2: $16 \div 2 = 8$ $\Rightarrow$ \textbf{Expression 8}
    \item \textbf{Expression 3:} 8 is even $\Rightarrow$ add 2: $8+2 = 10 \leq 15$
    $\Rightarrow$ \textbf{Expression 10}
    \item \textbf{BCD section:} Variant 14 is even $\Rightarrow$ use the \emph{even}
    column numbers: 21, 25, 36, 98, 125, 156
\end{itemize}

% ============================================
\newpage
\section{KEY INSTRUCTIONS USED}
% ============================================

\begin{table}[H]
\centering
\caption{NASM instructions used in this laboratory}
\small
\begin{tabular}{|l|p{9.5cm}|}
\hline
\textbf{Instruction} & \textbf{Description} \\
\hline
\texttt{MOV dest, src} & Copies a value from source to destination register or memory \\
\hline
\texttt{ADD dest, src} & Adds source to destination; result stored in destination \\
\hline
\texttt{SUB dest, src} & Subtracts source from destination \\
\hline
\texttt{IMUL dest, src} & Signed integer multiplication \\
\hline
\texttt{CDQ} & Sign-extends EAX into the EDX:EAX pair; required before \texttt{IDIV} \\
\hline
\texttt{IDIV src} & Signed division of EDX:EAX by src; quotient $\to$ EAX, remainder $\to$ EDX \\
\hline
\texttt{DAA} & Decimal Adjust AL after Addition; corrects AL to a valid packed BCD byte
after a binary \texttt{ADD} (32-bit mode only) \\
\hline
\texttt{JNC label} & Jump if Carry Flag = 0 (no carry generated) \\
\hline
\texttt{XOR reg, reg} & Zeros a register efficiently ($r \oplus r = 0$) \\
\hline
\texttt{INT 0x80} & Linux 32-bit syscall interrupt \\
\hline
\texttt{SYSCALL} & Linux 64-bit syscall instruction \\
\hline
\end{tabular}
\end{table}

\subsection{The DAA Instruction}

\texttt{DAA} (Decimal Adjust AL after Addition) is a special x86 instruction that
corrects the binary result of an \texttt{ADD} on two packed BCD bytes.  Its internal
logic is:

\begin{enumerate}
    \item If the low nibble of AL exceeds 9, or the Auxiliary Carry flag (AF) is set:
    $\texttt{AL} \mathrel{+}= 6$
    \item If AL exceeds \texttt{0x99}, or the Carry flag (CF) is set:
    $\texttt{AL} \mathrel{+}= \texttt{0x60}$; set CF
\end{enumerate}

\noindent\textbf{Important:} \texttt{DAA} is only valid in 32-bit (IA-32) mode.  Intel
removed it from the x86-64 instruction set.  Code using \texttt{DAA} must therefore be
assembled with \texttt{nasm -f elf32} and linked with \texttt{ld -m elf\_i386}.

% ============================================
\newpage
\section{TASK 1 --- REGISTER EXPRESSION}
% ============================================

\subsection{Formula and Values}

\[
A = (A + B) - (C + D)
\]

Test values: $A = 10$, $B = 5$, $C = 3$, $D = 2$.

Expected result: $(10 + 5) - (3 + 2) = 15 - 5 = \mathbf{10}$.

\subsection{Register Usage}

\begin{table}[H]
\centering
\caption{Register usage --- Task 1}
\begin{tabular}{|l|l|}
\hline
\textbf{Register} & \textbf{Contents} \\
\hline
EAX & Holds $A+B = 15$, then the final result $= 10$ \\
\hline
EBX & Holds $C+D = 5$ \\
\hline
\end{tabular}
\end{table}

\subsection{Source Code}

\begin{lstlisting}[caption={task1\_expression.asm --- A = (A+B) - (C+D)}]
; Task 1: A = (A + B) - (C + D)
; Works in SASM (64-bit mode, entry point: main)

section .data
    A      dd 10
    B      dd  5
    C      dd  3
    D      dd  2
    result dd  0

section .text
    global main

main:
    mov eax, [A]       ; eax = A
    add eax, [B]       ; eax = A + B

    mov ebx, [C]       ; ebx = C
    add ebx, [D]       ; ebx = C + D

    sub eax, ebx       ; eax = (A+B) - (C+D)

    mov [result], eax  ; store result

    xor eax, eax       ; return 0
    ret
\end{lstlisting}

\subsection{Step-by-Step Trace}

\begin{table}[H]
\centering
\caption{Execution trace --- Task 1}
\small
\begin{tabular}{|l|l|l|}
\hline
\textbf{Instruction} & \textbf{EAX} & \textbf{EBX} \\
\hline
\texttt{mov eax, [A]}  & 10 & --- \\
\hline
\texttt{add eax, [B]}  & 15 & --- \\
\hline
\texttt{mov ebx, [C]}  & 15 & 3 \\
\hline
\texttt{add ebx, [D]}  & 15 & 5 \\
\hline
\texttt{sub eax, ebx}  & \textbf{10} & 5 \\
\hline
\end{tabular}
\end{table}

% ============================================
\newpage
\section{TASK 2 --- ARITHMETIC EXPRESSIONS}
% ============================================

% -------------------------------------------
\subsection{Expression 14}
% -------------------------------------------

\subsubsection{Formula and Manual Verification}

\[
z = \frac{a \cdot c - b \cdot d}{f} + \frac{(-a+b) \cdot c}{d} \bigg/ h
\]

Values: $a=6,\ b=3,\ c=9,\ d=3,\ f=3,\ h=2$.

\[
z = \frac{6 \cdot 9 - 3 \cdot 3}{3} + \frac{(-6+3) \cdot 9}{3} \bigg/ 2
  = \frac{45}{3} + \frac{-27}{3} \bigg/ 2
  = \frac{15 + (-9)}{2}
  = \frac{6}{2} = \mathbf{3}
\]

\subsubsection{Register Usage}

\begin{table}[H]
\centering
\caption{Register usage --- Expression 14}
\small
\begin{tabular}{|l|p{9cm}|}
\hline
\textbf{Register} & \textbf{Contents at key steps} \\
\hline
EAX & $a \cdot c = 54 \to 54 - 9 = 45 \to 45/3 = 15 \to b-a=-3 \to -3 \cdot c=-27 \to -27/d=-9 \to -9+15=6 \to 6/2=3$ \\
\hline
EBX & $b \cdot d = 9$ \\
\hline
ECX & First term: $(a \cdot c - b \cdot d)/f = 15$ (saved across part 2) \\
\hline
EDX & Sign-extension half of dividend (set by \texttt{cdq}) \\
\hline
\end{tabular}
\end{table}

\subsubsection{Source Code}

\begin{lstlisting}[caption={task2\_expr14.asm --- Expression 14}]
; Task 2 - Expression 14: z = ((a*c - b*d)/f + (-a+b)*c/d) / h
; Values: a=6, b=3, c=9, d=3, f=3, h=2  =>  z = 3
; Works in SASM (64-bit mode, entry point: main)

section .data
    a      dd  6
    b      dd  3
    c      dd  9
    d      dd  3
    f      dd  3
    h      dd  2
    result dd  0

section .text
    global main

main:
    ; --- part 1: (a*c - b*d) / f ---
    mov  eax, [a]
    imul eax, dword [c]    ; eax = a*c = 54
    mov  ebx, [b]
    imul ebx, dword [d]    ; ebx = b*d = 9
    sub  eax, ebx          ; eax = 45
    cdq
    idiv dword [f]         ; eax = 45/3 = 15
    mov  ecx, eax          ; save first term

    ; --- part 2: (-a+b)*c / d ---
    mov  eax, [b]
    sub  eax, [a]          ; eax = b-a = -3
    imul eax, dword [c]    ; eax = -3*9 = -27
    cdq
    idiv dword [d]         ; eax = -27/3 = -9

    ; --- part 3: add terms, divide by h ---
    add  eax, ecx          ; eax = -9+15 = 6
    cdq
    idiv dword [h]         ; eax = 6/2 = 3  => z

    mov  [result], eax

    xor  eax, eax
    ret
\end{lstlisting}

% -------------------------------------------
\newpage
\subsection{Expression 8}
% -------------------------------------------

\subsubsection{Formula and Manual Verification}

\[
z = \frac{(-a+b+c+1)^{3}}{(a - b \cdot c + d)^{2}}
\]

Values: $a=2,\ b=3,\ c=1,\ d=2$.

\[
\text{num base} = -2+3+1+1 = 3 \quad\Rightarrow\quad 3^3 = 27
\]
\[
\text{den base} = 2 - 3 \cdot 1 + 2 = 1 \quad\Rightarrow\quad 1^2 = 1
\]
\[
z = 27 / 1 = \mathbf{27}
\]

\subsubsection{Register Usage}

\begin{table}[H]
\centering
\caption{Register usage --- Expression 8}
\small
\begin{tabular}{|l|p{9cm}|}
\hline
\textbf{Register} & \textbf{Contents at key steps} \\
\hline
EAX & base$=3 \to$ base$^2=9 \to$ base$^3=27$ (numerator) $\to z=27$ \\
\hline
EBX & base $= 3$ (saved for repeated multiplication) \\
\hline
ECX & numerator $= 27$ \\
\hline
EDX & denominator base $= 1 \to$ squared $= 1$ \\
\hline
\end{tabular}
\end{table}

\subsubsection{Source Code}

\begin{lstlisting}[caption={task2\_expr8.asm --- Expression 8}]
; Task 2 - Expression 8: z = (-a+b+c+1)^3 / (a - b*c + d)^2
; Values: a=2, b=3, c=1, d=2  =>  z = 27
; nasm -f elf64 task2_expr8.asm -o task2_expr8.o && ld task2_expr8.o -o task2_expr8
; Works in SASM (64-bit mode, entry point: main)

section .data
    a      dd  2
    b      dd  3
    c      dd  1
    d      dd  2
    result dd  0

section .text
    global main

main:
    ; --- numerator base: -a+b+c+1 ---
    mov  eax, [b]
    sub  eax, [a]          ; eax = b-a = 1
    add  eax, [c]          ; eax = 2
    add  eax, 1            ; eax = 3

    ; --- base^3 ---
    mov  ebx, eax          ; save base
    imul eax, ebx          ; eax = 9
    imul eax, ebx          ; eax = 27  (numerator)
    mov  ecx, eax

    ; --- denominator base: a - b*c + d ---
    mov  eax, [b]
    imul eax, dword [c]    ; eax = b*c = 3
    mov  edx, [a]
    sub  edx, eax          ; edx = a - b*c = -1
    add  edx, [d]          ; edx = 1

    ; --- base^2 ---
    imul edx, edx          ; edx = 1  (denominator)

    ; --- z = numerator / denominator ---
    mov  edi, edx          ; save denominator before cdq
    mov  eax, ecx          ; eax = 27 (numerator)
    cdq
    idiv edi               ; eax = 27/1 = 27  => z

    mov  [result], eax

    xor  eax, eax
    ret
\end{lstlisting}

% -------------------------------------------
\newpage
\subsection{Expression 10}
% -------------------------------------------

\subsubsection{Formula and Manual Verification}

\[
z = \frac{-a^{2} + b^{2}}{a^{2} - b^{2} - 5}
\]

Values: $a=5,\ b=3$.

\[
\text{numerator}   = -25 + 9 = -16
\qquad
\text{denominator} = 25 - 9 - 5 = 11
\]
\[
z = -16 \div 11 = \mathbf{-1} \quad\text{(integer division, truncates toward zero)}
\]

\subsubsection{Register Usage}

\begin{table}[H]
\centering
\caption{Register usage --- Expression 10}
\small
\begin{tabular}{|l|p{9cm}|}
\hline
\textbf{Register} & \textbf{Contents at key steps} \\
\hline
EAX & $a^2=25$ (temp) $\to$ numerator $=-16 \to z=-1$ \\
\hline
EBX & $b^2 = 9$ \\
\hline
ECX & numerator $= -16$ \\
\hline
EDX & denominator $= 11$ \\
\hline
ESI & $a^2 = 25$ (saved for reuse in both numerator and denominator) \\
\hline
\end{tabular}
\end{table}

\subsubsection{Source Code}

\begin{lstlisting}[caption={task2\_expr10.asm --- Expression 10}]
; Task 2 - Expression 10: z = (-a^2 + b^2) / (a^2 - b^2 - 5)
; Values: a=5, b=3  =>  z = -1
; nasm -f elf64 task2_expr10.asm -o task2_expr10.o && ld task2_expr10.o -o task2_expr10
; Works in SASM (64-bit mode, entry point: main)

section .data
    a      dd  5
    b      dd  3
    result dd  0

section .text
    global main

main:
    ; --- a^2 ---
    mov  eax, [a]
    imul eax, eax          ; eax = 25
    mov  esi, eax          ; save a^2

    ; --- b^2 ---
    mov  ebx, [b]
    imul ebx, ebx          ; ebx = 9

    ; --- numerator: -a^2 + b^2 ---
    mov  ecx, ebx
    sub  ecx, esi          ; ecx = 9-25 = -16

    ; --- denominator: a^2 - b^2 - 5 ---
    mov  edx, esi
    sub  edx, ebx          ; edx = 16
    sub  edx, 5            ; edx = 11

    ; --- z = numerator / denominator ---
    mov  eax, ecx
    cdq
    idiv edx               ; eax = -16/11 = -1  => z

    mov  [result], eax

    xor  eax, eax
    ret
\end{lstlisting}

% ============================================
\newpage
\section{TASK 3 --- BCD MANUAL CONVERSIONS (VARIANT 14, EVEN COLUMN)}
% ============================================

Even variant $\Rightarrow$ numbers from the even column: \textbf{21, 25, 36, 98, 125, 156}.

% -------------------------------------------
\subsection{Decimal to 8421 BCD}
% -------------------------------------------

\begin{table}[H]
\centering
\caption{Decimal to BCD conversion}
\begin{tabular}{|c|l|l|}
\hline
\textbf{Decimal} & \textbf{Digit breakdown} & \textbf{8421 BCD} \\
\hline
21  & 2 = 0010,\ 1 = 0001                       & \texttt{0010 0001} \\
\hline
25  & 2 = 0010,\ 5 = 0101                       & \texttt{0010 0101} \\
\hline
36  & 3 = 0011,\ 6 = 0110                       & \texttt{0011 0110} \\
\hline
98  & 9 = 1001,\ 8 = 1000                       & \texttt{1001 1000} \\
\hline
125 & 1 = 0001,\ 2 = 0010,\ 5 = 0101           & \texttt{0001 0010 0101} \\
\hline
156 & 1 = 0001,\ 5 = 0101,\ 6 = 0110           & \texttt{0001 0101 0110} \\
\hline
\end{tabular}
\end{table}

% -------------------------------------------
\subsection{Straight Binary vs.\ BCD --- Bit Count}
% -------------------------------------------

\begin{table}[H]
\centering
\caption{Binary bits vs.\ BCD bits required}
\small
\begin{tabular}{|c|l|c|l|c|c|}
\hline
\textbf{Decimal} & \textbf{Binary} & \textbf{Bits} & \textbf{BCD} & \textbf{BCD bits} & \textbf{Extra} \\
\hline
21  & 10101     & 5 & 0010 0001        & 8  & +3 \\
\hline
25  & 11001     & 5 & 0010 0101        & 8  & +3 \\
\hline
36  & 100100    & 6 & 0011 0110        & 8  & +2 \\
\hline
98  & 1100010   & 7 & 1001 1000        & 8  & +1 \\
\hline
125 & 1111101   & 7 & 0001 0010 0101   & 12 & +5 \\
\hline
156 & 10011100  & 8 & 0001 0101 0110   & 12 & +4 \\
\hline
\end{tabular}
\end{table}

BCD always requires more bits than straight binary because six of the sixteen possible
nibble values (1010--1111) are unused, wasting storage.

% -------------------------------------------
\subsection{BCD to Decimal Conversion}
% -------------------------------------------

\begin{table}[H]
\centering
\caption{BCD to decimal}
\begin{tabular}{|l|l|c|}
\hline
\textbf{BCD string} & \textbf{Split into nibbles} & \textbf{Decimal} \\
\hline
\texttt{001000110111}   & 0010 $|$ 0011 $|$ 0111 $\to$ 2, 3, 7 & 237 \\
\hline
\texttt{010000100001}   & 0100 $|$ 0010 $|$ 0001 $\to$ 4, 2, 1 & 421 \\
\hline
\texttt{100000000000}   & 1000 $|$ 0000 $|$ 0000 $\to$ 8, 0, 0 & 800 \\
\hline
\end{tabular}
\end{table}

% -------------------------------------------
\subsection{Adding BCD Binary Strings}
% -------------------------------------------

\textbf{(f) \ 01100100 + 00110011} \quad (decimal: $64 + 33 = 97$)

\begin{verbatim}
  01100100    (6, 4)
+ 00110011    (3, 3)
-----------
  10010111   -> nibbles: 1001 | 0111 = 9, 7  -> 97  (no correction needed)
\end{verbatim}

\textbf{(h) \ 10000101 + 00010011} \quad (decimal: $85 + 13 = 98$)

\begin{verbatim}
  10000101    (8, 5)
+ 00010011    (1, 3)
-----------
  10011000   -> nibbles: 1001 | 1000 = 9, 8  -> 98  (both <= 9, valid)
\end{verbatim}

\textbf{(e) \ 00100101 + 00100111} \quad (decimal: $25 + 27 = 52$)

\begin{verbatim}
Ones nibble:  0101 + 0111 = 1100 (12 > 9)  -> INVALID
              add 6: 1100 + 0110 = 0010, carry = 1
Tens nibble:  0010 + 0010 + 1(carry) = 0101 = 5  (valid)
Result: 0101 0010  ->  52  (correct)
\end{verbatim}

\textbf{(g) \ 10011000 + 10010111} \quad (decimal: $98 + 97 = 195$)

\begin{verbatim}
Ones nibble:  1000 + 0111 = 1111 (15 > 9)  -> add 6 -> 0101, carry=1
Tens nibble:  1001 + 1001 + 1   = 10011    -> add 6 -> 1001, carry=1
Hundreds:     0001
Result: 0001 1001 0101  ->  195  (correct)
\end{verbatim}

% -------------------------------------------
\newpage
\subsection{Decimal Pairs to BCD and Addition}
% -------------------------------------------

\textbf{a) \ 421 + 975 = 1396}

\begin{table}[H]
\centering
\small
\begin{tabular}{|l|l|}
\hline
421 in BCD & \texttt{0100 0010 0001} \\
\hline
975 in BCD & \texttt{1001 0111 0101} \\
\hline
Ones: $0001+0101$  & $= 0110 = 6 \leq 9$ (valid) \\
\hline
Tens: $0010+0111$  & $= 1001 = 9 \leq 9$ (valid) \\
\hline
Hundreds: $0100+1001$ & $= 1101 = 13 > 9$ $\to$ add 6: carry=1, digit=3 \\
\hline
Thousands: $0+0+1$ & $= 0001$ \\
\hline
\textbf{Result BCD} & \texttt{0001 0011 1001 0110} $\to$ \textbf{1396} \\
\hline
\end{tabular}
\end{table}

\textbf{b) \ 183 + 195 = 378}

\begin{table}[H]
\centering
\small
\begin{tabular}{|l|l|}
\hline
183 in BCD & \texttt{0001 1000 0011} \\
\hline
195 in BCD & \texttt{0001 1001 0101} \\
\hline
Ones: $0011+0101$  & $= 1000 = 8 \leq 9$ (valid) \\
\hline
Tens: $1000+1001$  & $= 10001 > 9$ $\to$ add 6 $\to$ 0111, carry=1 \\
\hline
Hundreds: $0001+0001+1$ & $= 0011 = 3$ (valid) \\
\hline
\textbf{Result BCD} & \texttt{0011 0111 1000} $\to$ \textbf{378} \\
\hline
\end{tabular}
\end{table}

% -------------------------------------------
\subsection{BCD Subtraction via 9's Complement}
% -------------------------------------------

\subsubsection{Method}
To compute $A - B$ in BCD:
\begin{enumerate}
    \item Find the 9's complement of $B$: subtract each decimal digit from 9.
    \item Add $A$ and the 9's complement of $B$ using BCD addition rules.
    \item If a carry is generated (end-around carry), add it to the result.
    That carry signals a positive result.
    \item If no carry, the result is negative; take the 9's complement of the sum.
\end{enumerate}

\subsubsection{a) \ 974 - 632 = 342}

9's complement of $632$: $9-6=3,\ 9-3=6,\ 9-2=7$ $\Rightarrow$ \textbf{367}

\begin{verbatim}
  974 + 367 (BCD)
  Ones:     0100 + 0111 = 1011 > 9  -> add 6 -> 0001, carry=1
  Tens:     0111 + 0110 + 1 = 1110 > 9  -> add 6 -> 0100, carry=1
  Hundreds: 1001 + 0011 + 1 = 1101 > 9  -> add 6 -> 0011, carry=1 (end-around)
  Sum = 0341 + end-around carry 1 = 342  ->  POSITIVE  (correct)
\end{verbatim}

\subsubsection{b) \ 453 - 983 = -530}

9's complement of $983$: $9-9=0,\ 9-8=1,\ 9-3=6$ $\Rightarrow$ \textbf{016}

\begin{verbatim}
  453 + 016 (BCD)
  Ones:     0011 + 0110 = 1001 = 9  (valid)
  Tens:     0101 + 0001 = 0110 = 6  (valid)
  Hundreds: 0100 + 0000 = 0100 = 4  (valid)
  Sum = 469, NO carry  ->  NEGATIVE
  9's complement of 469: 9-4=5, 9-6=3, 9-9=0  ->  530
  Result: -530  (correct)
\end{verbatim}

\subsubsection{Large numbers: X = 72345, Y = -24341, \ X + Y}

Compute $72345 - 24341$ using 5-digit 9's complement.

9's complement of $24341$: $99999 - 24341 = 75658$

\begin{verbatim}
  72345 + 75658 (BCD):
  Ones:         5+8=13 > 9  -> add 6 -> 3, carry=1
  Tens:         4+5+1=10 > 9  -> add 6 -> 0, carry=1
  Hundreds:     3+6+1=10 > 9  -> add 6 -> 0, carry=1
  Thousands:    2+5+1=8  (valid)
  Ten-thousands: 7+7=14 > 9  -> add 6 -> 4, carry=1  (end-around)
  Sum = 48003 + 1 = 48004  ->  POSITIVE
  72345 + (-24341) = 48004  (correct)
\end{verbatim}

% ============================================
\newpage
\section{TASK 4 --- BCD ADDER CIRCUIT IN LOGISIM}
% ============================================

\subsection{Circuit Description}

A single-digit BCD adder takes two 4-bit BCD digits ($A_{3:0}$ and $B_{3:0}$) and a
carry-in ($C_{in}$), and produces a 4-bit BCD sum digit plus a carry-out ($C_{out}$).

The circuit consists of two main stages:

\begin{enumerate}
    \item \textbf{Binary addition:} A standard 4-bit binary full-adder (FA) adds
    $A_{3:0}$, $B_{3:0}$, and $C_{in}$.  The raw 4-bit result $S_{3:0}$ and the
    carry-out $C_4$ are produced.
    \item \textbf{Correction logic:} The correction (add 6 = 0110) is needed when the
    raw sum is invalid BCD, i.e.\:
    \[
        \text{CORRECT} = C_4 \;+\; (S_3 \cdot S_2) \;+\; (S_3 \cdot S_1)
    \]
    If CORRECT = 1, a second 4-bit adder adds 0110 to $S_{3:0}$.  Any carry from this
    second addition becomes the BCD $C_{out}$.
\end{enumerate}

\begin{table}[H]
\centering
\caption{BCD adder --- Logisim component list}
\begin{tabular}{|l|p{9cm}|}
\hline
\textbf{Component} & \textbf{Purpose} \\
\hline
4-bit binary adder (×2) & First adds the BCD digits; second optionally adds the correction 0110 \\
\hline
3-input OR gate & Computes CORRECT $= C_4 + S_3 S_2 + S_3 S_1$ \\
\hline
2-input AND gate (×2) & Computes $S_3 \cdot S_2$ and $S_3 \cdot S_1$ \\
\hline
4-bit MUX / AND gates & Selects 0110 as second operand of second adder when CORRECT = 1 \\
\hline
OR gate for $C_{out}$ & Combines carry from second adder with CORRECT signal \\
\hline
\end{tabular}
\end{table}

\subsection{Truth Table for Correction Signal}

\begin{table}[H]
\centering
\caption{When correction is applied (selected cases)}
\small
\begin{tabular}{|c|c|c|c|c|c|}
\hline
$C_4$ & $S_3$ & $S_2$ & $S_1$ & CORRECT & Action \\
\hline
0 & 0 & x & x & 0 & No correction \\
\hline
0 & 1 & 0 & 0 & 0 & No correction (sum = 8 or 9) \\
\hline
0 & 1 & 0 & 1 & 1 & Add 6 (sum = 10 or 11) \\
\hline
0 & 1 & 1 & x & 1 & Add 6 (sum = 12--15) \\
\hline
1 & x & x & x & 1 & Add 6 (carry means sum $\geq$ 16) \\
\hline
\end{tabular}
\end{table}

% ============================================
\newpage
\section{TASK 5 --- ASSEMBLY: BCD ADDITION}
% ============================================

\subsection{Description}

The program adds two packed 8-bit BCD numbers and stores the two-byte result (low digit
byte and carry byte) in sequential memory locations.  Because the \texttt{DAA} instruction
is unavailable in 64-bit mode, the code is assembled as a 32-bit ELF binary.

\begin{itemize}
    \renewcommand{\labelitemi}{--}
    \item \textbf{Build:} \texttt{nasm -f elf32 task5\_bcd\_add.asm -o task5\_bcd\_add.o}
    \item \textbf{Link:} \texttt{ld -m elf\_i386 task5\_bcd\_add.o -o task5\_bcd\_add}
    \item \textbf{Note:} On a 64-bit system you may need \texttt{sudo apt install gcc-multilib}
    (Debian/Ubuntu) or \texttt{sudo pacman -S lib32-glibc} (Arch Linux) for 32-bit support.
\end{itemize}

\textbf{Example:} BCD $0x47$ $(=47)$ $+$ BCD $0x38$ $(=38)$

\begin{enumerate}
    \item Binary sum: $0x47 + 0x38 = 0x7F$
    \item \texttt{DAA}: lower nibble \texttt{F} = 15 $> 9$ $\Rightarrow$ AL $\mathrel{+}= 6$:
    $0x7F + 0x06 = 0x85$
    \item Upper nibble $= 8 \leq 9$ --- no second correction; CF $= 0$
    \item Result: AL $= 0x85$ (BCD 85), carry $= 0$
\end{enumerate}

\subsection{Source Code}

\begin{lstlisting}[caption={task5\_bcd\_add.asm --- Add two 8-bit BCD numbers}]
; Task 5: Add two 8-bit packed BCD numbers
; DAA only works in 32-bit mode!
; In SASM: Settings -> Assembler -> check "Use 32-bit program"
;
; Example: 0x47 + 0x38 = 0x85  (47 + 38 = 85),  carry = 0

section .data
    bcd1  db 0x47          ; first BCD number  (47)
    bcd2  db 0x38          ; second BCD number (38)

section .bss
    result    resb 1       ; BCD result byte
    carry_out resb 1       ; carry byte (0x00 or 0x01)

section .text
    global main

main:
    mov  al, [bcd1]        ; al = 0x47
    add  al, [bcd2]        ; al = 0x7F  (raw binary, not valid BCD yet)
    daa                    ; fix al to valid BCD: 0x7F -> 0x85
                           ; DAA adds 6 to any nibble that exceeds 9

    mov  [result], al      ; store result

    jnc  no_carry
    mov  byte [carry_out], 0x01   ; result > 99, carry to next digit
    jmp  done

no_carry:
    mov  byte [carry_out], 0x00

done:
    xor  eax, eax
    ret
\end{lstlisting}

\subsection{Memory and Register Analysis}

\begin{table}[H]
\centering
\caption{Memory state before and after execution}
\begin{tabular}{|l|c|c|p{5cm}|}
\hline
\textbf{Location} & \textbf{Before} & \textbf{After} & \textbf{Meaning} \\
\hline
\texttt{[bcd1]}      & 0x47 & 0x47 & BCD 47 (input, unchanged) \\
\hline
\texttt{[bcd2]}      & 0x38 & 0x38 & BCD 38 (input, unchanged) \\
\hline
\texttt{[result]}    & 0x00 & 0x85 & BCD 85 (packed: tens=8, ones=5) \\
\hline
\texttt{[carry\_out]} & 0x00 & 0x00 & No carry (result $\leq$ 99) \\
\hline
\end{tabular}
\end{table}

\begin{table}[H]
\centering
\caption{Register trace during execution}
\small
\begin{tabular}{|l|l|}
\hline
\textbf{After instruction} & \textbf{AL value and flags} \\
\hline
\texttt{mov al, [bcd1]}  & AL = 0x47 \\
\hline
\texttt{add al, [bcd2]}  & AL = 0x7F \ (AF=1 because lower nibble 7+8=15 $>$ 9) \\
\hline
\texttt{daa} --- lower check  & Nibble F $= 15 > 9$ $\to$ AL $+= 6$: $0x7F + 0x06 = 0x85$ \\
\hline
\texttt{daa} --- upper check  & Nibble 8 $\leq 9$ $\to$ no correction; CF $= 0$ \\
\hline
\texttt{mov [result], al}  & \texttt{[result]} $= 0x85$ (BCD 85) \\
\hline
\texttt{jnc no\_carry}  & CF $= 0$ $\to$ jump taken \\
\hline
\texttt{mov byte [carry\_out], 0x00}  & \texttt{[carry\_out]} $= 0x00$ \\
\hline
\end{tabular}
\end{table}

% ============================================
\newpage
\section{DEBUGGING}
% ============================================

Debugging Assembly code requires careful attention to register and memory state at every
step.  The following issues were encountered during development:

\begin{table}[H]
\centering
\caption{Debugging issues and solutions}
\small
\begin{tabular}{|p{4cm}|p{5cm}|p{5cm}|}
\hline
\textbf{Issue} & \textbf{Cause} & \textbf{Solution} \\
\hline
\texttt{IDIV} produced wrong quotient or fault & \texttt{CDQ} was missing before
\texttt{IDIV}; EDX contained leftover data from a previous multiply & Add \texttt{cdq}
immediately before every \texttt{idiv} to sign-extend EAX into EDX:EAX \\
\hline
Expression 8: \texttt{idiv edx} crashed (division by zero) & After \texttt{cdq}, EDX is
overwritten; the denominator stored in EDX was destroyed & Compute and save the
denominator in EDX \emph{before} calling \texttt{cdq}/\texttt{idiv} \\
\hline
Task 5 assembled for 64-bit (elf64) & \texttt{DAA} is not supported in x86-64 mode;
NASM gave an ``instruction not supported in 64-bit mode'' error & Switched to
\texttt{nasm -f elf32} and \texttt{ld -m elf\_i386} \\
\hline
\texttt{IMUL ebx, dword [d]} used with a memory operand & Two-operand \texttt{IMUL}
requires the destination to be a register; using a memory operand gives an assemble error
& Load memory value into a register first, then \texttt{IMUL} \\
\hline
\end{tabular}
\end{table}

\subsection{Debugging Approach}

All programs were debugged using \textbf{GDB} with the following workflow:

\begin{enumerate}
    \item Assemble with debug symbols: add the \texttt{-g} flag to NASM.
    \item Start GDB: \texttt{gdb ./task2\_expr14}
    \item Set breakpoint at the division instruction: \texttt{break *\_start+N}
    \item Run: \texttt{run}
    \item Inspect registers: \texttt{info registers eax ebx ecx edx}
    \item Single-step: \texttt{stepi} (one instruction at a time)
    \item Check memory: \texttt{x/4xb \&result}
\end{enumerate}

This approach made it easy to verify intermediate values (e.g.\ confirming EAX $= 45$
before the first \texttt{idiv} in Expression 14) and to catch the missing \texttt{cdq}
bug early.

% ============================================
\newpage
\section{CONCLUSION}
% ============================================

This laboratory covered two complementary topics in low-level computing: register-based
arithmetic in NASM Assembly and the BCD (Binary Coded Decimal) number system.

Through Task 1 and Task 2, the mechanics of signed integer arithmetic using \texttt{IMUL},
\texttt{IDIV}, and \texttt{CDQ} were practised.  A key lesson was that \texttt{CDQ} must
always precede \texttt{IDIV}, and that intermediate results must be saved in registers
such as ECX or ESI before they are overwritten by subsequent operations.

The BCD section (Task 3) demonstrated why BCD is used in contexts where exact decimal
representation matters (e.g.\ financial software, seven-segment displays): each decimal
digit is stored independently, so no rounding errors or base-conversion artefacts occur.
The manual addition and 9's complement subtraction exercises illustrated how the six
invalid states (1010--1111) necessitate the $+6$ correction step.

Task 4 showed how the correction logic can be implemented in hardware using two 4-bit
adders and a small combinational circuit.  Task 5 demonstrated the software equivalent
via the \texttt{DAA} instruction, which encapsulates exactly the same correction rules in
a single opcode.  The important architectural note is that \texttt{DAA} was removed in
x86-64, requiring 32-bit mode for any code that depends on it.

Overall, the laboratory strengthened understanding of how computers represent and
manipulate numbers at the hardware level, and highlighted the gap between the
abstractions provided by high-level languages and the actual operations performed by the
processor.

\end{document}
